\chapter{Training Configuration Details}
\label{app:yolox_config}

\section{QR and Barcode Detection Model}

\begin{lstlisting}[caption={YOLOX experiment file for QR and barcode detection}, label={lst:yolox_barcode}]
"""
YOLOX experiment configuration for QR and barcode detection.

This configuration defines the training setup used for the custom
barcode detection model. The model is based on the YOLOX-S architecture
and trained on a merged dataset containing both real and synthetic data.

The configuration includes dataset paths, training parameters, and
augmentation settings used during training.
"""

from yolox.exp import Exp as MyExp
from yolox.data import COCODataset, TrainTransform, ValTransform
import os


class Exp(MyExp):
    def __init__(self):
        super(Exp, self).__init__()

        # Experiment name
        self.exp_name = "barcode_merged_yolox_s"

        # --------------------------------------------------
        # Model configuration (YOLOX-S)
        # --------------------------------------------------
        self.depth = 0.33
        self.width = 0.50
        self.num_classes = 1  # Single class: barcode / QR / visual code

        # --------------------------------------------------
        # Dataset configuration
        # --------------------------------------------------
        self.data_dir = os.path.abspath(
            os.path.join(os.path.dirname(__file__), "../../../../barcodeV2/dataset_merged")
        )

        # Annotation files (COCO format)
        self.train_ann = "instances_train2017.json"
        self.val_ann = "instances_test2017.json"  # Used for evaluation

        # Image directories
        self.train_name = "train2017"
        self.val_name = "test2017"

        # --------------------------------------------------
        # Training configuration
        # --------------------------------------------------
        self.max_epoch = 200
        self.data_num_workers = 4

        # Higher resolution improves detection of small barcodes
        # but increases GPU memory usage
        self.input_size = (960, 960)
        self.test_size = (960, 960)

        self.print_interval = 20
        self.eval_interval = 1

        # --------------------------------------------------
        # Data augmentation configuration
        # --------------------------------------------------
        # Mosaic and MixUp disabled due to already diverse dataset
        self.mosaic_prob = 0.0
        self.mixup_prob = 0.0

        # Moderate online augmentations
        self.hsv_prob = 0.5        # Color jitter
        self.flip_prob = 0.5       # Horizontal flip

        # Geometric transformations (kept small to preserve structure)
        self.degrees = 5.0
        self.translate = 0.05
        self.shear = 1.0

    def get_dataset(self, cache=False, cache_type=None):
        """
        Returns the training dataset using COCO format.
        """
        return COCODataset(
            data_dir=self.data_dir,
            json_file=self.train_ann,
            name=self.train_name,
            img_size=self.input_size,
            preproc=TrainTransform(
                max_labels=120,
                flip_prob=self.flip_prob,
                hsv_prob=self.hsv_prob,
            ),
            cache=cache,
            cache_type=cache_type,
        )

    def get_eval_dataset(self, testdev=False, legacy=False, **kwargs):
        """
        Returns the evaluation dataset.
        """
        return COCODataset(
            data_dir=self.data_dir,
            json_file=self.val_ann,
            name=self.val_name,
            img_size=self.test_size,
            preproc=ValTransform(legacy=legacy),
        )
\end{lstlisting}

\section{License Plate Detection Model}

\begin{lstlisting}[caption={YOLOX experiment file for license plate detection}, label={lst:yolox_licenseplate}]
"""
YOLOX experiment configuration for license plate detection.

This configuration defines the training and evaluation setup for the
license plate detection model. The model is based on YOLOX-S and trained
on a custom dataset with COCO-style annotations.

Compared to the barcode model, this configuration uses more aggressive
data augmentation (e.g., mosaic) to improve generalization.
"""

#!/usr/bin/env python3
# -*- coding:utf-8 -*-

import os
import torch
from yolox.exp import Exp as MyExp
from yolox.data import COCODataset, ValTransform
from yolox.evaluators import COCOEvaluator


class Exp(MyExp):
    def __init__(self):
        super().__init__()

        # Experiment name
        self.exp_name = "yolox_s_license_plate"

        # --------------------------------------------------
        # Model configuration (YOLOX-S)
        # --------------------------------------------------
        self.depth = 0.33
        self.width = 0.50
        self.num_classes = 1  # Single class: license plate

        # --------------------------------------------------
        # Dataset configuration
        # --------------------------------------------------
        self.data_dir = os.path.abspath(
            os.path.join(os.path.dirname(__file__), "../../../../license_plate_detection/dataset")
        )

        self.train_ann = "train.json"
        self.val_ann = "test.json"  # Used only for evaluation

        # --------------------------------------------------
        # Training configuration
        # --------------------------------------------------
        self.input_size = (960, 960)
        self.test_size = (960, 960)

        self.max_epoch = 200
        self.no_aug_epochs = 15

        # Learning rate schedule
        self.basic_lr_per_img = 0.01 / 64.0
        self.min_lr_ratio = 0.05
        self.warmup_epochs = 5

        self.data_num_workers = 2
        self.print_interval = 20
        self.eval_interval = 1

        # --------------------------------------------------
        # Data augmentation configuration
        # --------------------------------------------------
        # Mosaic enabled to improve generalization
        self.mosaic_scale = (0.5, 1.5)
        self.mosaic_prob = 0.8

        # MixUp disabled
        self.mixup_prob = 0.0
        self.enable_mixup = False

        # Standard augmentations
        self.hsv_prob = 0.5
        self.flip_prob = 0.5

        # Geometric transformations
        self.degrees = 5.0
        self.translate = 0.05
        self.shear = 1.0

        # Multi-scale training
        self.random_size = (14, 24)

        # --------------------------------------------------
        # Evaluation settings
        # --------------------------------------------------
        self.test_conf = 0.01
        self.nmsthre = 0.65

        self.output_dir = "./YOLOX_outputs"

    def get_eval_loader(self, batch_size, is_distributed, testdev=False, legacy=False):
        """
        Creates the evaluation data loader.
        """
        image_folder = "test2017" if self.val_ann == "test.json" else "val2017"

        valdataset = COCODataset(
            data_dir=self.data_dir,
            json_file=self.val_ann,
            name=image_folder,
            img_size=self.test_size,
            preproc=ValTransform(legacy=legacy),
        )

        if is_distributed:
            batch_size = batch_size // torch.distributed.get_world_size()
            sampler = torch.utils.data.distributed.DistributedSampler(
                valdataset, shuffle=False
            )
        else:
            sampler = torch.utils.data.SequentialSampler(valdataset)

        return torch.utils.data.DataLoader(
            valdataset,
            batch_size=batch_size,
            sampler=sampler,
            num_workers=self.data_num_workers,
            pin_memory=True,
        )

    def get_evaluator(self, batch_size, is_distributed, testdev=False, legacy=False):
        """
        Returns the COCO evaluator used for model evaluation.
        """
        return COCOEvaluator(
            dataloader=self.get_eval_loader(batch_size, is_distributed, testdev, legacy),
            img_size=self.test_size,
            confthre=self.test_conf,
            nmsthre=self.nmsthre,
            num_classes=self.num_classes,
            testdev=testdev,
        )
\end{lstlisting}